% ============================================================
% discussion.tex
% ============================================================

This chapter synthesizes our empirical results, connects them
to the broader literature and policy landscape, and identifies
the most important directions for extending this work. We
organize the discussion around four substantive themes: the
PSLRA as a procedural filter, judicial geography as a
structural determinant of outcomes, the role of case-level
characteristics, and the methodological lessons from
comparing semi-parametric and nonparametric frameworks.

\section{Summary of the Research Question and Approach}

This thesis set out to answer a deceptively simple question:
what determines whether a federal securities class action
settles, is dismissed, and how quickly? Answering it
correctly requires confronting two statistical challenges
that prior work has not simultaneously addressed. First,
settlement and dismissal are competing risks: they are
mutually exclusive terminal events, and studying settlement
duration conditional on settlement occurring recovers a
distorted picture of litigation dynamics. Second, the
dominant covariate of interest---PSLRA regime---has a
time-varying effect whose procedural mechanisms operate
through the motion-to-dismiss phase, implying a hazard ratio
that changes over litigation duration.

We address both challenges using a competing-risks survival
framework applied to 12,968 federal securities class actions
from the FJC Integrated Database. The framework combines
Aalen-Johansen nonparametric CIF estimation, cause-specific
Cox models, Fine-Gray subdistribution models, and a Random
Survival Forest benchmark, all evaluated on a held-out test
set. To our knowledge, this is the first application of a
full competing-risks framework to population-level
administrative data on federal securities litigation.

\section{Policy Implications of the PSLRA Effect}

Our results provide the most precise empirical
characterization to date of how the Private Securities
Litigation Reform Act of 1995 reshaped the litigation
landscape. The three-phase temporal decomposition is the
single most policy-relevant finding in this thesis. In the
first year, post-PSLRA cases face more than double the
dismissal hazard of pre-PSLRA cases (HR $= 2.18$). This is
the PSLRA's heightened pleading standard and automatic
discovery stay operating as designed: the ``strong
inference'' of scienter requirement raises the bar for
surviving a motion to dismiss, and the discovery stay
prevents plaintiffs from using early discovery to supplement
insufficiently pleaded complaints.

The null effect in years one to two (HR $= 0.922$, $p =
0.444$) is equally important. Cases that survive the
first-year gauntlet face no differential dismissal risk
compared to pre-PSLRA survivors in the following year. This
pattern is consistent with the PSLRA acting as a one-time
filter at the pleading stage: once a case has demonstrated
sufficient pleading adequacy to survive initial review, the
reform no longer systematically disadvantages plaintiffs.
The re-emergence of a positive PSLRA effect beyond two years
(HR $= 1.55$) may reflect differences in summary judgment
practice or a selection effect among long-duration
post-PSLRA survivors.

From a policy standpoint, our findings support the
interpretation that the PSLRA achieved its primary
legislative objective of filtering meritless early-stage
claims. The first-year dismissal hazard amplification is
large and precisely estimated; the subsequent null effect
suggests that surviving cases proceed on approximately equal
footing with pre-PSLRA cases. Whether this filtering is
appropriately calibrated---that is, whether some meritorious
cases are caught in the same net---requires merging our
dataset with settlement amount data and merit proxies, and
remains an open empirical question.

The robustness analysis provides an important bound on the
settlement effect. Under the most liberal coding scheme
(Scheme~B), the PSLRA settlement HR rises from 0.378 to
0.696, still representing a 30\% reduction in settlement
hazard. The true effect lies within the interval
$[0.378, 0.696]$ depending on the fraction of Code~12 cases
representing private settlements. Even at the upper bound,
the PSLRA substantially suppressed settlement probability,
consistent with defendants' reduced incentive to settle
cases they could credibly dismiss at the pleading stage.

\section{Circuit Heterogeneity and Judicial Geography}

The dominance of circuit geography in our results---confirmed
by both the parametric Cox models and the nonparametric RSF
variable importance rankings---raises fundamental questions
about the role of judicial institutions in shaping litigation
outcomes. The Second Circuit's anomalous position is the
most striking geographic finding. Cases filed there are
anywhere from 5 to 17 times less likely to settle at any
given duration than cases in other major circuits, yet they
have the highest cumulative dismissal incidence. This
combination defines a litigation environment distinctively
hostile to monetary recovery for plaintiffs' counsel.

Several mechanisms could explain this pattern. The Southern
District of New York handles a disproportionate share of the
most complex and high-stakes securities cases, and defendants
in these cases are typically represented by elite defense
counsel with substantial expertise, which may reduce early
settlement probability. SDNY judges also have extensive
experience with securities class actions and may apply
demanding standards at class certification, reducing the
expected value of cases to plaintiffs. The PSLRA $\times$
Circuit interaction analysis adds a temporal dimension:
the PSLRA's suppression of settlement hazard was largest in
the Second Circuit, consistent with the hypothesis that the
Act's heightened pleading standard interacted with SDNY's
already-demanding judicial environment.

The Eleventh Circuit's anomalous interaction for dismissal
(interaction HR $= 0.426$) and the Ninth Circuit's
amplification (HR $= 1.052$) point to meaningful variation
in how district courts applied the PSLRA's ``strong
inference'' scienter standard. Prior work by Johnson et al.\
\citep{johnson2001} and Choi \citep{choi2007} has documented
circuit-level variation in PSLRA application; our
competing-risks framework provides the first quantitative
estimates of how this variation manifests in the joint
distribution of settlement and dismissal timing.

The policy implication is sobering. If the PSLRA's filtering
effects are substantially larger in the Second Circuit, then
plaintiffs filing in New York face a qualitatively different
litigation environment than those filing elsewhere. This
raises questions about regulatory forum shopping, the
appropriate geographic scope of securities law reform, and
whether the PSLRA inadvertently created
jurisdiction-specific barriers to recovery.

\section{Structural Case Characteristics: MDL, Origin,
and Statutory Basis}

The MDL finding is methodologically novel. MDL
consolidation suppresses both cause-specific hazards by
approximately 73--74\%, indicating that MDL cases resolve
at dramatically lower instantaneous rates regardless of
their eventual outcome. The natural interpretation is that
consolidation introduces coordination costs and procedural
delays---bellwether trials, coordinated discovery
schedules, global resolution frameworks---that slow
individual case resolution. The Fine-Gray subdistribution
analysis complicates this picture in an important way:
despite the dramatically lower instantaneous hazards, MDL
cases do not have lower cumulative settlement probability
(SHR $= 0.963$). Because bilateral suppression of both
hazards keeps large numbers of MDL cases in the risk set
throughout the study period, cumulative settlement
probability accumulates toward the same long-run level as
non-MDL cases. This finding would be completely invisible
to a single-outcome analysis.

Case origin results suggest that removed cases face lower
dismissal hazard (HR $= 0.783$) and directionally higher
settlement hazard, consistent with the interpretation that
cases initially filed in state court involve claims where
settlement is a more natural resolution than federal
dismissal. Statutory basis results are subtler: Section~11
cases have a modestly higher dismissal hazard than
Section~10(b) cases (HR $= 1.121$), which may reflect
tracing and standing requirements in IPO and secondary
offering contexts that create early dismissal opportunities
defendants exploit aggressively.

\section{The Semi-Parametric Advantage (and Computational Limits)}

The comparison between our semi-parametric and nonparametric 
models yields a stark divergence in out-of-sample discrimination: 
the Cox model substantially outperforms the current RSF implementation 
across all outcomes and horizons. The Cox settlement C-index of 0.735 
versus the RSF's 0.549 represents a gap of nearly 0.19 concordance units.

While it is tempting to interpret this as definitive evidence that 
securities litigation outcomes are strictly dominated by linear main 
effects (circuit identity and PSLRA regime) without complex interactions, 
such a conclusion is premature. The current RSF performance (barely 
exceeding random chance) is highly characteristic of an under-tuned 
ensemble. With heavily right-censored competing-risks data, 200 trees 
and a restricted $m_{\text{try}}$ grid are insufficient to stabilize 
the cumulative hazard estimates or reliably map complex recursive 
partitions. 

Therefore, the proper interpretation at this interim stage is that 
the Cox model's partial likelihood efficiently captures the massive, 
first-order institutional shocks (the PSLRA) and geographic disparities 
(Circuit effects). Whether a fully scaled, computationally intensive 
RSF can uncover subtle, nonlinear interactions among secondary case 
characteristics remains an open empirical question for the final 
computation.

\section{Limitations}

The primary limitations of our modeling framework center on unobserved 
heterogeneity. Our hazard ratios represent population-averaged effects 
conditional only on observed administrative covariates. However, prior 
literature identifies case-specific merit signals---such as the magnitude 
of stock price drops, the presence of SEC investigations, or accounting 
restatements---as strong predictors of outcomes \citep{cox2008, johnson2001}. 
Because these metrics are absent from our dataset, our models cannot isolate 
how much of the PSLRA's early dismissal effect was driven by the dismissal 
of genuinely meritless claims versus structurally complex but valid claims. 


Furthermore, litigation outcomes are highly dependent on the specific 
judge assigned to the case and the caliber of the law firms involved 
\citep{wang2022}. Without consistent judge identifiers or attorney data, 
we cannot implement frailty models to account for this clustering. 
Consequently, unobserved judicial or counsel quality may confound some 
of our observed circuit or origin effects. Finally, our reliance on the 
Grambsch-Therneau global tests indicates that proportional hazards 
assumptions are violated in our extended models; our current estimates 
must therefore be strictly interpreted as average effects over the study 
period rather than constant instantaneous risks.

\section{Future Research Directions}

Several extensions would substantially strengthen this
analysis. CourtListener integration would provide
title-based class action validation, docket activity
complexity proxies, and attorney identifiers, enabling
examination of law firm quality effects in a competing-risks
duration framework \citep{wang2022}. Merging with the
Stanford Securities Class Action Clearinghouse would provide
merit proxies---stock price drop magnitudes, accounting
restatement indicators, parallel SEC action flags---that
would sharpen interpretation of PSLRA filtering effects.
Full time-varying coefficient models for circuit, origin,
MDL, and statutory basis are needed to address the
documented proportional hazards assumption violations.
Frailty models incorporating judge-level random effects,
if consistent judge identifiers can be recovered from
CourtListener, could account for unobserved judge-level
heterogeneity. Settlement amount analysis combined with
our competing-risks duration framework would provide a more
complete picture of the economic consequences of PSLRA
and circuit effects. Finally, testing for post-2018
structural breaks---including the Supreme Court's decision
in \textit{Cyan, Inc.\ v.\ Beaver County Employees
Retirement Fund}---is a natural extension of the piecewise
hazard approach developed here.

\section{Conclusion}

Federal securities class actions are among the most economically 
consequential civil cases litigated in the United States. 
 
Yet, their duration and outcome dynamics have historically been evaluated 
through incomplete statistical lenses that ignore the fundamentally 
competing nature of case resolution. By applying a full competing-risks 
survival framework to population-level administrative data, this thesis 
corrects that methodological blind spot and reveals a litigation landscape 
defined by profound structural asymmetries.

We have demonstrated that the PSLRA transformed securities litigation 
not as a permanent dampener on plaintiff success, but as a highly specific, 
first-year procedural filter. We have shown that multidistrict litigation 
consolidation dramatically slows the pace of justice without systematically 
altering the ultimate probability of settlement. Most significantly, we 
have quantified how circuit geography dictates case trajectories, revealing 
a fragmented federal system where the jurisdiction of filing is as 
determinative of a case's fate as the substantive law itself. 

As policymakers continue to debate the balance between deterring corporate 
fraud and curbing frivolous litigation, the findings presented here provide 
a rigorous empirical foundation. Evaluating the health of the securities 
litigation system requires looking beyond simple settlement averages; 
it demands an accounting of the competing risks, procedural delays, and 
geographic lotteries that define the modern federal docket.